% Research continuation of the pinned ProveIt snapshot
% 4c173c06cc32c9cea554b39be1837ad2ae897fc1, 2026-10-10.
% This is an insertion fragment, not a stand-alone document.
% Prior sources: chapters/05-real-euler.tex, 05-universal-euler.tex;
% incoming ProveIt_Polylogarithms_Research_2026-10-10/
% polylogarithms_uniform_bounds_20261010/sections/euler.tex.

\section{The optimal late Euler constant: eventual axis reduction}
\label{lateaxis:sec:global}

The incoming uniform-bounds report~\cite{incoming-bounds} proves $C_N\to1$ and an explicit
lower asymptotic $\liminf N^p(C_N-1)\ge c$.  It leaves the matching
upper asymptotic open.  We prove the stronger fact that, after an
explicit (deliberately conservative) index, the exact global supremum
is the maximum on the zero outer-order axis.  This also identifies
the extremizing inner order and a second asymptotic scale.

Write
\[
 W_N(x)=\frac{(1-x)^N}{1+x},\qquad
 F_N(y)=W_N(e^{-2y}),\qquad
 J_N(s)=\mathbb E F_N(T_s),\qquad D_N(s)=1-J_N(s),
\]
where $T_s$ is a rate-one gamma variable of shape $s>0$ and $T_0=0$.
Thus $J_N(0)=0$.  Define
\[
 R_N(a,b)=2^N(E_N(a,b)-g_{a,b}),\qquad
 C_N=\sup_{a,b>0}R_N(a,b),\qquad
 A_N=\max_{b>0}R_N(0,b).
\]
The exact representation already established in the manuscript is
\begin{equation}\label{lateaxis:eq:kernel}
 R_N(a,b)=\mathbb E\frac{F_N(T_a+U_b)-F_N(T_a)}{1-e^{-U_b}},
\end{equation}
with independent rate-one gamma variables.  Its positive gamma-series
form, valid also when $b\le1$, is
\begin{equation}\label{lateaxis:eq:series}
 R_N(a,b)=\sum_{r\ge1}r^{-b}
 \mathbb E\{F_N(T_a+U_b/r)-F_N(T_a)\}.
\end{equation}
For $b>1$ this gives the useful upper bound
\begin{equation}\label{lateaxis:eq:zetaJ}
 R_N(a,b)\le\zeta(b)J_N(a+b).
\end{equation}
Indeed $F_N$ is increasing, so every first expectation in
\eqref{lateaxis:eq:series} is at most $J_N(a+b)$, while the subtracted
expectations are nonnegative.  The known strict estimate
$R_N(a,b)<1$ for $a\ge1$, $b>0$ is the first case of the manuscript's
Theorem~\texttt{rw:thm:euler}.  For clarity, it can also be recovered from the compensated density.
At $a=1$,
\[
 K_{1,b}(T)=-\frac{T^b}{\Gamma(b+1)}+h_b(T).
\]
Each of the two positive terms has $e^{-T}\,dT$ integral one;
for $h_b$, interchange its positive integrals and use
$(1-e^{-v})/(e^v-1)=e^{-v}$. Their strict overlap makes the positive
and negative Jordan masses strictly less than one. For $a>1$,
convolution with $T^{a-2}/\Gamma(a-1)$ in the $T$ variable preserves
zero total weighted mass and contracts the $L^1(e^{-T}dT)$ norm.
Since $0\le W_N\le1$, the exact remainder integral is bounded above
by the negative Jordan mass, and hence is strictly less than one.
This gives the needed bound directly for every real $a\ge1$.

\subsection{A gamma-transform shape principle}

We shall use the following elementary observation twice.  Let
$h(y)>0$ be continuous, increase strictly to one maximum, and then
decrease strictly, with distinct finite endpoint limits $h(0)$ and
$h(\infty)$.  Suppose each horizontal level between the smaller
endpoint limit and the maximum crosses $h$ either once or twice,
as its shape prescribes.  Consider
\[
 H(s)=\frac{r^s}{\Gamma(s)}\int_0^\infty
                  h(y)y^{s-1}e^{-ry}\,dy,\qquad s>0.
\]
For the bounded functions used below, every derivative in $s$ exists.
At a stationary point $s_0$, set $\lambda=H(s_0)$ and
$f=h-\lambda$.  Since $H(s_0)=\lambda$ and $H'(s_0)=0$,
\[
 \int f(y)y^{s_0-1}e^{-ry}\,dy
 =\int f(y)\log y\,y^{s_0-1}e^{-ry}\,dy=0.
\]
If $f$ has just one sign change at $y_1$, multiplication by
$\log y-\log y_1$ produces an integrand of one strict sign,
a contradiction.  If $f$ has signs $-,+,-$ and zeros $y_1<y_2$,
then $-(\log y-\log y_1)(\log y-\log y_2)$ has the same signs.
Its integral against $f(y)y^{s_0-1}e^{-ry}\,dy$ is strictly positive.
The two vanishing integrals therefore imply $H''(s_0)<0$.
Consequently every stationary point is a nondegenerate maximum.
There is at most one, and the minimum of $H$ on any compact interval
is attained at an endpoint.  Existence of one value above both
endpoint limits implies a unique global maximum.

\begin{lemma}[Quasiconcavity of a tilted defect]\label{lateaxis:lem:defect}
For every integer $N\ge2$, the function
\[
 Q_N(s)=2^sD_N(s),\qquad s>0,
\]
has no interior minimum; in particular, its minimum on every compact
interval is attained at an endpoint.
\end{lemma}
\begin{proof}
It is the gamma transform of shape $s$ and rate two of
\[
 \psi_N(y)=e^y\{1-W_N(e^{-2y})\}.
\]
This bounded positive function has endpoint limits one and zero.
We check its strict one-maximum shape.  Put $x=e^{-2y}$ and
\[
 B_N(x)=\frac{(1-x)^{N-1}
                 \{1+2(N+1)x+(2N-3)x^2\}}{(1+x)^2}.
\]
Direct differentiation gives
\[
 \psi_N'(y)=e^y\{1-B_N(x)\}.
\]
The numerator of the logarithmic derivative of $B_N$ after
multiplication by the positive denominator
$(1-x)(1+x)\{1+2(N+1)x+(2N-3)x^2\}$ is
\begin{align*}
 P_N(x)={}&N+1-(2N^2+N+5)x\notag\\
          &-(4N^2-3N-7)x^2-(2N-3)(N-1)x^3.
\end{align*}
For $N\ge2$ every nonconstant coefficient is negative, so $P_N$
decreases strictly and has exactly one zero in $(0,1)$.
Since $B_N(0)=1$, $B_N(1)=0$, it first increases and then decreases,
and crosses the level one exactly once in $(0,1)$.  It follows that
$\psi_N$ first increases and then decreases.  The preceding
gamma-transform principle, at rate two, proves the claim.
\end{proof}

\subsection{Excluding bounded and moderately growing inner orders}

\begin{lemma}[An explicit region below one]\label{lateaxis:lem:exclude}
Let $N$ be an integer with $N\ge e^{24}$ and write $\ell=\log N$.
Then
\[
 R_N(a,b)<1\qquad(0\le a\le1,\quad 0<b\le2\ell).
\]
\end{lemma}
\begin{proof}
First suppose $0<b\le3$.  Direct differentiation yields, for $N\ge2$,
\[
 0\le F_N'(y)\le4Nx(1-x)^{N-1}\le4,
 \qquad x=e^{-2y}.
\]
Couple independent gamma variables so that $T_a\le T_1$ and
$U_b\le U_3$, and put $V=T_1+U_3$.  Since
$u/(1-e^{-u})\le1+u$, the exact kernel is at most
\[
 (1+U_3)\sup_{0\le y\le V}F_N'(y).
\]
On $V\le10$, set $z=Nx\ge Ne^{-20}\ge e^4>50$.  Then
\[
 F_N'(y)\le4z e^{-z/2}\le200e^{-25}.
\]
On $V>10$ use $F_N'\le4$.  Now $V$ has shape four, and the
conditional mean of $U_3$ given $V$ is $3V/4$.  The elementary
integer-shape gamma tails give
\begin{align*}
 R_N(a,b)
 &\le800e^{-25}+4\{\mathbb P(T_4>10)
                         +3\mathbb P(T_5>10)\}\\
 &=800e^{-25}+\frac{25928}{3}e^{-10}<1.
\end{align*}
All constants in this estimate are explicit; for instance $e>27/10$
suffices to verify the last inequality.

Next suppose $3\le b\le2\ell$.  By
\eqref{lateaxis:eq:zetaJ} and the increase of $J_N$,
\[
 R_N(a,b)\le\zeta(b)\{1-D_N(b+1)\}.
\]
We prove $2^bD_N(b+1)>4$ on this entire interval by checking its
endpoints and using Lemma~\ref{lateaxis:lem:defect}.

At $b=3$, Chebyshev's inequality gives
$\mathbb P(T_4\le8)\ge3/4$.  Since
$W_N(x)\le e^{-Nx}$,
\[
 D_N(4)\ge\frac34\{1-e^{-Ne^{-16}}\}
        \ge\frac34\{1-e^{-e^8}\}>\frac12.
\]
At $b=2\ell$, use
\[
 1-W_N(x)\ge1-(1-x)^N\ge\frac{Nx}{1+Nx}
                       \ge\frac{Nx}{2}\quad(0\le x\le N^{-1}).
\]
Thus, for $s=2\ell+1$ and a gamma variable $Z_s$ of shape $s$
and rate three,
\[
 D_N(s)\ge\frac N2\,3^{-s}\,\mathbb P(Z_s\ge\ell/2).
\]
Its mean is $s/3$ and its variance is $s/9$, so
\[
 \mathbb P(Z_s<\ell/2)
 \le\frac{4(2\ell+1)}{(\ell+2)^2}\le\frac12
 \qquad(\ell\ge24).
\]
Writing $L=\log(3/2)<5/12$, we conclude that
\[
 2^{2\ell}D_N(2\ell+1)
 \ge\frac{N^{1-2L}}{12}
 >\frac{e^4}{12}>4.
\]
The bound $L<5/12$ follows from the trapezoidal bound for
$\int_1^{3/2}dt/t$.  Lemma~\ref{lateaxis:lem:defect}, applied on
$4\le s\le2\ell+1$, proves $D_N(b+1)>4\,2^{-b}$ throughout.
Finally,
\[
 \zeta(b)-1\le2^{-b}\left(1+\frac2{b-1}\right)\le2\,2^{-b}.
\]
Hence $\zeta(b)(1-D_N(b+1))<1$, as required.
\end{proof}

\subsection{Outer-order monotonicity in the remaining region}

For a bounded Lipschitz function $f$, gamma addition gives the exact
shape-generator formula
\begin{equation}\label{lateaxis:eq:generator}
 \frac{d}{da}\mathbb E f(T_a)
 =\int_0^\infty\frac{e^{-u}}u
             \mathbb E\{f(T_a+u)-f(T_a)\}\,du.
\end{equation}
It includes the right derivative at $a=0$.  Indeed, divide the
gamma-addition identity at $a+h$ by $h$ and use
$h\Gamma(h)\to1$.  A Lipschitz bound at zero and boundedness at
infinity justify dominated convergence.

\begin{lemma}[An explicit strictly decreasing outer-order range]
\label{lateaxis:lem:outer}
If $N\ge e^{24}$, $b\ge2\log N$, and $0\le a\le1$, then
\[
 \partial_a R_N(a,b)<0,
\]
with the one-sided interpretation at $a=0$.
\end{lemma}
\begin{proof}
Put $t_0=\frac12\log N$.  In \eqref{lateaxis:eq:generator} for
$J_N$, restrict to $T_a\le1$ and $u\in[t_0+1,t_0+2]$.
Gamma addition gives $\mathbb P(T_a\le1)\ge1-e^{-1}$.  For $N\ge16$,
\[
 F_N(T_a+u)-F_N(T_a)
 \ge F_N(t_0+1)-F_N(1)>\frac12.
\]
To check the last inequality, use
$1-W_N(x)\le(N+1)x$, $W_N(x)\le e^{-Nx}$, and $7<e^2<8$.
Consequently
\begin{equation}\label{lateaxis:eq:Jprime}
 J_N'(a)\ge
 \frac{(1-e^{-1})e^{-2}}{\sqrt N(\log N+4)}
 >\frac{1}{18\sqrt N(\log N+4)}.
\end{equation}

For $b>1$, rewrite the positive series as
\[
 R_N(a,b)=\zeta(b)\{1-J_N(a)\}-\mathcal D_N(a,b),
 \quad
 \mathcal D_N(a,b)=\sum_{r\ge1}r^{-b}
       \mathbb E\{1-F_N(T_a+U_b/r)\}.
\]
The derivative of $1-W_N(x)$ lies in $[0,N+1]$, and therefore
\[
 0\le F_N(y+u)-F_N(y)
 \le(N+1)e^{-2y}(1-e^{-2u}).
\]
Apply \eqref{lateaxis:eq:generator} to each summand.  Positivity and
the following summable bound justify interchanging derivative and sum:
\begin{align}
 0\le-\partial_a\mathcal D_N(a,b)
 &\le (N+1)\log3\,3^{-a}\sum_{k\ge3}k^{-b}.
 \label{lateaxis:eq:Dprime}
\end{align}
Here $\int_0^\infty(e^{-u}-e^{-3u})\,du/u=\log3$, and
$r^{-b}\mathbb E e^{-2U_b/r}=(r+2)^{-b}$.
Since $b\ge2\log N\ge3$,
\[
 \sum_{k\ge3}k^{-b}\le3^{-b}\left(1+\frac3{b-1}\right)
 \le3\,3^{-b}\le3N^{-2},
\]
where $3>e$ was used in the last step.  Thus the right side of
\eqref{lateaxis:eq:Dprime} is at most $12/N$.
By \eqref{lateaxis:eq:Jprime},
\[
 \partial_a R_N(a,b)
 <-\frac1{18\sqrt N(\log N+4)}+\frac{12}{N}<0.
\]
The final inequality follows from
$\sqrt N>216(\log N+4)$ for $\log N\ge24$.
\end{proof}

\begin{theorem}[Eventual exact reduction of the global supremum]
\label{lateaxis:thm:reduction}
For every integer $N\ge\lceil e^{24}\rceil$,
\[
 \boxed{\quad C_N=A_N=\max_{b>0}R_N(0,b).\quad}
\]
Every positive outer order gives a strict inequality
$R_N(a,b)<A_N$.  If the axis $a=0$ is admitted, the unique maximizing
pair is $(0,b_N)$, where $b_N$ is the unique maximum of the axis
function proved below.  The threshold is an explicit sufficient
condition, not an estimate of the smallest index for which the
conclusion holds.
\end{theorem}
\begin{proof}
The axis maximum is greater than one.  This follows already from the
exact Taylor lower bound
\[
 R_N(0,b)\ge1+2^{-b}-N3^{-b}-N4^{-b}>1
\]
for all sufficiently large $b$ at fixed $N$.
For $a\ge1$ the previously proved strict bound is $R_N(a,b)<1$.
For $0\le a\le1$ and $b\le2\log N$ use
Lemma~\ref{lateaxis:lem:exclude}.  In the remaining range,
Lemma~\ref{lateaxis:lem:outer} gives
$R_N(a,b)<R_N(0,b)\le A_N$ when $a>0$.
Finally the bounded kernel in \eqref{lateaxis:eq:kernel} is continuous
under $a\downarrow0$, so the axis values are approached by positive
outer orders.  This proves equality of the suprema and all strict
claims.  Uniqueness on the axis follows from the next theorem.
\end{proof}

\section{The maximizing inner order and a fractional second scale}
\label{lateaxis:sec:asymptotic}

On the axis the exact integral is
\[
 A_N(b):=R_N(0,b)=\frac1{\Gamma(b)}\int_0^\infty
        y^{b-1}e^{-y}K_N(e^{-y})\,dy,
 \qquad K_N(u)=\frac{(1+u)(1-u^2)^{N-1}}{1+u^2}.
\]
In this section $A_N(b)$ denotes the function and $A_N$ without an
argument its maximum.

\begin{theorem}[A unique axis maximum at every index]
\label{lateaxis:thm:unique}
For each integer $N\ge2$, $A_N(b)$ has exactly one stationary point
$b_N>0$.  It is a nondegenerate global maximum, greater than one;
$A_N'$ is positive below $b_N$ and negative above it.
\end{theorem}
\begin{proof}
The sign of $K_N'(u)$ on $0<u<1$ is the sign of
\[
 1-(2N+1)u+u^2-(2N-3)u^3.
\]
This polynomial decreases strictly, from a positive value to a
negative value, because its derivative is at most $-(2N+1)+2<0$.
Thus $K_N(e^{-y})$ increases strictly from zero to one maximum,
then decreases strictly to one.  The gamma-transform principle
proves that all stationary points of $A_N$ are nondegenerate
maxima.  Moreover $A_N(b)\to0$ as $b\downarrow0$,
$A_N(b)\to1$ as $b\to\infty$, and the preceding Taylor lower
bound gives values greater than one.  A global maximum therefore
exists, and the shape principle makes it unique.
\end{proof}

\begin{corollary}[Monotone height and maximizing inner order]
\label{lateaxis:cor:monotone}
For every integer $N\ge2$,
\[
 A_{N+1}<A_N,\qquad b_{N+1}>b_N.
\]
In particular $C_N$ is strictly decreasing for
$N\ge\lceil e^{24}\rceil$.
\end{corollary}
\begin{proof}
The pointwise relation $K_{N+1}(u)=(1-u^2)K_N(u)<K_N(u)$
implies $A_N(b_{N+1})>A_{N+1}(b_{N+1})=A_{N+1}$,
and hence $A_N>A_{N+1}$.  For the locations, the ratio
$A_{N+1}(b)/A_N(b)$ is the expectation of
$1-e^{-2T}$ under the probability density proportional to
$K_N(e^{-T})T^{b-1}e^{-T}$.  Its derivative in $b$ is the strictly
positive covariance of $1-e^{-2T}$ and $\log T$ under that density.
Thus $A_{N+1}'(b_N)>0$, since $A_N'(b_N)=0$.
The unique-maximum theorem gives $b_{N+1}>b_N$.
The eventual reduction transfers the height inequality to $C_N$.
\end{proof}

Set
\begin{gather*}
 L=\log(3/2),\qquad B=L^{-1},\qquad
 p=B\log2,\qquad q=\frac{\log3}{\log2},\qquad
 d_0=B\log q,\\
 c=(1-q^{-1})q^{-p},\qquad
 \nu=B-\frac12,\qquad
 I=\frac12-B+B\log(2B),\\
 D=\sqrt{\frac{B}{2\pi}}\,\Gamma\left(\frac12-B\right)
                          (2B)^{-d_0},\qquad
 K=\frac{\log(2B)D}{(\log2)Lq^{-p}}.
\end{gather*}
The elementary bounds $2/5<L<5/12$ show $1<\nu<2$, so
$D>0$.  Also $p<I<2p-1$.

\begin{theorem}[The sharp global late asymptotic and maximizing order]
\label{lateaxis:thm:sharpasymptotic}
As $N\to\infty$, with $\ell=\log N$,
\begin{align}
 C_N=A_N
 &=1+cN^{-p}
     +D\frac{N^{-I}}{\sqrt\ell}
     +o\left(\frac{N^{-I}}{\sqrt\ell}\right),
 \label{lateaxis:eq:second}\\
 b_N
 &=\frac{\log(Nq)}{\log(3/2)}
   -K\frac{N^{-(I-p)}}{\sqrt\ell}
   +o\left(\frac{N^{-(I-p)}}{\sqrt\ell}\right).
 \label{lateaxis:eq:location}
\end{align}
In particular the incoming conjecture is settled:
\[
 \lim_{N\to\infty}N^p(C_N-1)=c.
\]
The numerical diagnostics are
\begin{align*}
 p&=1.70951129135\ldots,& c&=0.16794779654\ldots,\\
 I&=1.96959041447\ldots,& D&=1.56777020071\ldots.
\end{align*}
The second correction comes from a gamma-density transition near
$y=\frac12\log N$; it is not the full quadratic Taylor moment.
\end{theorem}
\begin{proof}
Put
\[
 P_N(s)=\frac{(1-s)^{N-1}}{1+s},\qquad
 h_N(t)=P_N(t/N)-1+t\quad(0\le t\le N).
\]
The elementary Taylor estimates
\[
 P_N(0)=1,\quad P_N'(0)=-N,\quad
 0\le P_N''(s)\le N^2-N+2
\]
give, for $N\ge2$,
\begin{equation}\label{lateaxis:eq:hbound}
 0\le h_N(t)\le\min(t,t^2).
\end{equation}
They also give the exact decomposition
\begin{equation}\label{lateaxis:eq:decompose}
 A_N(b)=1+2^{-b}-N3^{-b}-N4^{-b}+M_N(b),
\end{equation}
where
\[
 M_N(b)=\mathbb E\{(1+e^{-T_b})h_N(Ne^{-2T_b})\}\ge0.
\]
For $b=B\ell+d$, where $d$ ranges over a fixed compact interval,
center this last integral at $t_0=\ell/2$.  If
$p_b(t)=t^{b-1}e^{-t}/\Gamma(b)$, then
\begin{align*}
 \frac{M_N(b)}{p_b(t_0)}
 =\int_{-t_0}^\infty &e^{-y}(1+2y/\ell)^{b-1}
       (1+N^{-1/2}e^{-y})h_N(e^{-2y})\,dy.
\end{align*}
For each fixed $y$ the integrand converges to
\[
 e^{(2B-1)y}\{e^{-e^{-2y}}-1+e^{-2y}\}.
\]
This passage is dominated, uniformly for bounded $d$.  Indeed, choose
$\varepsilon>0$ so that $2<2B-1-\varepsilon<2B-1+\varepsilon<4$.
The inequality $\log(1+v)\le v$ and \eqref{lateaxis:eq:hbound}
give an integrable majorant proportional to
$e^{(2B-1-\varepsilon-2)y}$ on $y<0$ and
$e^{(2B-1+\varepsilon-4)y}$ on $y>0$.
The factor $1+N^{-1/2}e^{-y}$ is at most two on the integration
domain.  Extending by zero below $-t_0$ completes the domination.

With $t=e^{-2y}$ the limiting integral equals
\begin{equation}\label{lateaxis:eq:mellin}
 \frac12\int_0^\infty
       (e^{-t}-1+t)t^{-\nu-1}\,dt
 =\frac12\Gamma(-\nu).
\end{equation}
Both integrations by parts used in the last equality are convergent
because $1<\nu<2$; equivalently the value is
$\Gamma(2-\nu)/(2\nu(\nu-1))$.
Stirling's formula, uniformly for bounded $d$, gives
\[
 p_{B\ell+d}(\ell/2)
 =\sqrt{\frac{2B}{\pi}}\,(2B)^{-d}
       \frac{N^{-I}}{\sqrt\ell}\{1+O(\ell^{-1})\}.
\]
Consequently
\begin{equation}\label{lateaxis:eq:Mlimit}
 M_N(B\ell+d)=D(d)\frac{N^{-I}}{\sqrt\ell}\{1+o(1)\},
 \quad D(d)=\sqrt{\frac{B}{2\pi}}\,\Gamma(-\nu)(2B)^{-d}.
\end{equation}

Differentiating the gamma density inserts the score
$\log(t_0+y)-\psi(b)$, which tends to $-\log(2B)$.
The same domination is valid after multiplication by this score.
For $y\ge0$ its absolute value is bounded by $C(1+y)$.
For $y<0$, put $z=1+2y/\ell\in(0,1)$ and use
$(-\log z)z\le1/e$ to replace $z^{b-1}$ by $z^{b-2}$;
this only changes the exponent of the preceding majorant by
$O(\ell^{-1})$.  Hence, uniformly for bounded $d$,
\begin{equation}\label{lateaxis:eq:Mprime}
 M_N'(B\ell+d)=-\log(2B)D(d)
                       \frac{N^{-I}}{\sqrt\ell}\{1+o(1)\}.
\end{equation}

Let $f(d)=2^{-d}-3^{-d}$.  Since $B\log3-1=p$, the first two
nonconstant terms in \eqref{lateaxis:eq:decompose} equal
$N^{-p}f(d)$.  The term $N4^{-b}$ and its derivative have order
$N^{1-2p}=o(N^{-I}/\sqrt\ell)$.  The inequalities $p<I$ and
$I<2p-1$ follow, for example, from
\[
 I-p=B\log(2B/3)-B+\frac32>0
\]
(strict convexity of $x\log x-x+1$), together with
$12/5<B<5/2$, $\log5<5/3$, and $\log2>2/3$ for the second
inequality.  The latter crude bounds give $I<13/6$ and
$2p-1>11/5$.

Now $f'$ is positive below $d_0$, negative above it, and
\[
 f(d_0)=c,\qquad f''(d_0)=-(\log2)Lq^{-p}<0.
\]
For every fixed $\delta>0$, equations
\eqref{lateaxis:eq:decompose} and \eqref{lateaxis:eq:Mprime}
give opposite derivative signs at $B\ell+d_0\pm\delta$ for all
sufficiently large $N$.  The unique critical point therefore has
$b_N=B\ell+d_0+o(1)$.  Substituting it into the stationary-point
equation and using $f''(d_0)\ne0$ yields
\eqref{lateaxis:eq:location}.  Its displacement changes the maximum
value by $o(N^{-I}/\sqrt\ell)$, so
\eqref{lateaxis:eq:decompose} and \eqref{lateaxis:eq:Mlimit} give
\eqref{lateaxis:eq:second} for $A_N$.  The exact eventual reduction
Theorem~\ref{lateaxis:thm:reduction} transfers it to $C_N$.
\end{proof}

\begin{corollary}[Concentration of almost maximizing orders]
\label{lateaxis:cor:concentration}
Let $a_N,b_N^{\circ}>0$ satisfy
\[
 C_N-R_N(a_N,b_N^{\circ})=o(N^{-p}).
\]
Then
\[
 a_N=o\bigl(N^{1/2-p}\log N\bigr),\qquad
 b_N^{\circ}=\frac{\log(Nq)}{\log(3/2)}+o(1).
\]
Thus orders attaining the leading excess over one must approach the
outer-order axis at a quantified rate.
\end{corollary}
\begin{proof}
The sharp asymptotic implies $R_N(a_N,b_N^{\circ})>1$ eventually,
so the exclusion argument forces $a_N<1$ and $b_N^{\circ}>2\log N$.
The estimates in Lemma~\ref{lateaxis:lem:outer} in fact give
\[
 \partial_aR_N(a,b)\le-\frac1{36\sqrt N(\log N+4)}
\]
on that region, because $\sqrt N>432(\log N+4)$ for $\log N\ge24$.
Integrating in $a$ yields
\[
 C_N-R_N(a,b)\ge A_N-A_N(b)
                     +\frac a{36\sqrt N(\log N+4)}.
\]
Both terms on the right are nonnegative, proving the bound on $a_N$.
For each fixed $\varepsilon>0$, the axis values at
$B\log N+d_0\pm\varepsilon$ have leading coefficients
$f(d_0\pm\varepsilon)<c$.  Axis unimodality therefore excludes
every $b$ outside this interval from having
$A_N-A_N(b)=o(N^{-p})$.  Let $\varepsilon\downarrow0$ to obtain
the asserted location of $b_N^{\circ}$.
\end{proof}

\paragraph{Scope and next questions.}
The proof establishes an eventual theorem with a conservative explicit
threshold; it does not establish $C_N=A_N$ at every $N\ge2$.
Determining the smallest such index, or proving equality at all indices,
is a precise remaining question.  A second question is an effective
error bound in \eqref{lateaxis:eq:second}; dominated convergence gives
the stated asymptotic but no small numerical error estimate at ordinary
truncation indices.  Expanding the centered gamma density suggests an
inverse-logarithm series multiplying $N^{-I}/\sqrt{\log N}$, with
coefficients given by derivatives of $\Gamma(-\nu)$, before still
smaller powers of $N$.  That longer expansion is not claimed here.
